% !PRINTED 681
\weberpage{681}
%<<CONT
schuss der Zahl $x$ über die grösste in $x$ enthaltene ganze Zahl
bedeutet, so ergiebt sich Folgendes:

Ist $\xi_1, \xi_2, \ldots, \xi_{\nu-1}$ das Exponentensystem einer Einheit, und
$m$ eine ganze Zahl, so ist auch

\weq{6}{(m\,\xi_1), (m\,\xi_2), \ldots, (m\,\xi_{\nu-1})}

das Exponentensystem einer Einheit.

Nun sind die Grössen $(m\,\xi_i)$, wenn sie nicht Null sind,
positive echte Brüche, und nach 1) muss es sich also, wenn die
Reihe der ganzen Zahlen hinlänglich weit fortgesetzt wird, ereignen,
dass für zwei verschiedene Werthe $m'$, $m''$ von $m$ die
Zahlenreihe (6) übereinstimmt. Da hiernach $m'\xi_i$, $m''\xi_i$ denselben
Ueberschuss über eine ganze Zahl haben, so ist $(m' - m'')\,\xi_i$
selbst eine ganze Zahl, und es giebt also eine ganze rationale
Zahl $m$, die höchstens gleich der Anzahl der Exponentensysteme
reducirter Einheiten ist, von der Eigenschaft, dass

\[
m\,\xi_1, m\,\xi_2, \ldots, m\,\xi_{\nu-1}
\]

ganze rationale Zahlen werden.

Diese Zahl $m$ kann sich ändern, wenn die Einheit $\varepsilon$ geändert
wird. Da aber alle möglichen Werthe dieser Zahl unter einer
endlichen Grenze liegen, so giebt es eine endliche Zahl, in der
alle diese Werthe von $m$ enthalten sind und die selbst für $m$
genommen werden kann. Es giebt daher eine gewisse positive
ganze Zahl $m$, die als Nenner aller der rationalen Zahlen genommen
 werden kann, die als Exponenten von Einheiten auftreten
 können. Damit ist der Satz 6. bewiesen.

\begin{center}
\S.\,187.\\[4pt]
\gesperrt{Fundamentalsysteme von Einheiten.}
\end{center}

Wir wählen jetzt an Stelle des Systemes unabhängiger Einheiten
 $\varepsilon_i$ ein anderes, dessen conjugirte Logarithmen

\weq{1}{\lambda_{i,1}, \lambda_{i,2}, \ldots, \lambda_{i,\nu} \qquad i = 1, 2, \ldots, \nu-1}

sein mögen. Nach \S.\,186, (5) ist

\weq{2}{\begin{aligned}
&\lambda_{i,k} = \xi_{1,i}\, l_{1,k} + \xi_{2,i}\, l_{2,k} + \cdots + \xi_{\nu-1,i}\, l_{\nu-1,k} \\
&i = 1, 2, \ldots, \nu-1; \quad k = 1, 2, \ldots, \nu,
\end{aligned}}

wenn $\xi_{1,i}, \xi_{2,i}, \ldots, \xi_{\nu-1,i}$ die Exponenten einer der neuen Einheiten
 in Bezug auf das System $\varepsilon_1, \varepsilon_2, \ldots, \varepsilon_{\nu-1}$ bedeuten, die
wir als rationale Zahlen mit dem gemeinsamen Nenner $m$ annehmen
 können.
